%%%%%%%%%%%%%%%%%%%%%%%%%%%%%%%%%%%%%%%%%%%%%%%%%%%%%%%%%%%%
% 14.4 MATHEMATICAL CONSTANTS
% The Composition of Advanced Mathematics
%%%%%%%%%%%%%%%%%%%%%%%%%%%%%%%%%%%%%%%%%%%%%%%%%%%%%%%%%%%%

\section{Mathematical Constants}
\label{sec:mathematical-constants}

\prerequisite{
Basic algebra, geometry, trigonometry, calculus, complex numbers,
probability, and mathematical notation.
}

\learningobjectives{
By the end of this reference section, the reader should be able to:
\begin{itemize}
    \item recognize several major mathematical constants;
    \item distinguish exact definitions from decimal approximations;
    \item identify the circle constant \(\pi\);
    \item identify Euler's number \(e\);
    \item understand the imaginary unit \(i\);
    \item recognize the golden ratio \(\phi\);
    \item recognize the Euler--Mascheroni constant \(\gamma\);
    \item understand common appearances of constants in analysis;
    \item recognize constants arising from special functions and number theory;
    \item distinguish universal mathematical constants from arbitrary
          parameters;
    \item recognize constants appearing in geometry;
    \item recognize constants appearing in probability and statistics;
    \item recognize constants appearing in asymptotic analysis;
    \item understand why decimal approximations should not replace exact forms
          unnecessarily;
    \item recognize that many constants are irrational or transcendental;
    \item use constants consistently in symbolic computation and numerical
          approximation.
\end{itemize}
}

%%%%%%%%%%%%%%%%%%%%%%%%%%%%%%%%%%%%%%%%%%%%%%%%%%%%%%%%%%%%
\subsection{What Is a Mathematical Constant?}
\label{subsec:what-is-mathematical-constant}
%%%%%%%%%%%%%%%%%%%%%%%%%%%%%%%%%%%%%%%%%%%%%%%%%%%%%%%%%%%%

A mathematical constant is a fixed numerical value that arises naturally in
mathematical definitions, theorems, limits, geometry, analysis, probability,
or other structures.

Examples include

\[
\boxed{
\pi,
\qquad
e,
\qquad
\phi,
\qquad
\gamma.
}
\]

Unlike a variable, a constant does not change within its mathematical
definition.

%%%%%%%%%%%%%%%%%%%%%%%%%%%%%%%%%%%%%%%%%%%%%%%%%%%%%%%%%%%%
\subsection{Exact Values and Decimal Approximations}
\label{subsec:exact-versus-decimal-constants}
%%%%%%%%%%%%%%%%%%%%%%%%%%%%%%%%%%%%%%%%%%%%%%%%%%%%%%%%%%%%

Many important constants have exact symbolic definitions but nonterminating
decimal expansions.

For example,

\[
\boxed{
\pi
\approx
3.141592653589793.
}
\]

The symbol \(\pi\) is exact.

The decimal

\[
3.141592653589793
\]

is only an approximation.

Whenever possible, exact symbolic forms should be preserved during
derivations.

%%%%%%%%%%%%%%%%%%%%%%%%%%%%%%%%%%%%%%%%%%%%%%%%%%%%%%%%%%%%
\subsection{Quick Constant Reference}
\label{subsec:quick-constant-reference}
%%%%%%%%%%%%%%%%%%%%%%%%%%%%%%%%%%%%%%%%%%%%%%%%%%%%%%%%%%%%

\begin{center}
\begin{tabular}{c|c|l}
\textbf{Constant}
&
\textbf{Approximate Value}
&
\textbf{Common Meaning}
\\ \hline

\(\pi\)
&
\(3.141592653589793\)
&
circle constant
\\

\(e\)
&
\(2.718281828459045\)
&
natural exponential base
\\

\(\phi\)
&
\(1.618033988749895\)
&
golden ratio
\\

\(\gamma\)
&
\(0.577215664901533\)
&
Euler--Mascheroni constant
\\

\(\sqrt{2}\)
&
\(1.414213562373095\)
&
diagonal ratio of a unit square
\\

\(\sqrt{3}\)
&
\(1.732050807568877\)
&
common geometric radical
\\

\(\ln 2\)
&
\(0.693147180559945\)
&
natural logarithm of \(2\)
\\

\(\zeta(2)\)
&
\(1.644934066848226\)
&
equals \(\pi^2/6\)
\\

\(G\)
&
\(0.915965594177219\)
&
Catalan's constant
\\

\(\delta\)
&
\(4.669201609102990\)
&
Feigenbaum constant for period doubling
\end{tabular}
\end{center}

%%%%%%%%%%%%%%%%%%%%%%%%%%%%%%%%%%%%%%%%%%%%%%%%%%%%%%%%%%%%
\subsection{The Circle Constant \(\pi\)}
\label{subsec:pi-constant-reference}
%%%%%%%%%%%%%%%%%%%%%%%%%%%%%%%%%%%%%%%%%%%%%%%%%%%%%%%%%%%%

The constant \(\pi\) is defined by the ratio

\[
\boxed{
\pi
=
\frac{
\text{circumference of a circle}
}{
\text{diameter of the circle}
}.
}
\]

For every Euclidean circle, this ratio is the same.

Its decimal expansion begins

\[
\boxed{
\pi
=
3.14159265358979323846\ldots
}
\]

%%%%%%%%%%%%%%%%%%%%%%%%%%%%%%%%%%%%%%%%%%%%%%%%%%%%%%%%%%%%
\subsection{Basic Circle Formulas Using \(\pi\)}
\label{subsec:pi-circle-formulas}
%%%%%%%%%%%%%%%%%%%%%%%%%%%%%%%%%%%%%%%%%%%%%%%%%%%%%%%%%%%%

For a circle of radius \(r\),

\[
\boxed{
C
=
2\pi r
}
\]

is the circumference, and

\[
\boxed{
A
=
\pi r^2
}
\]

is the area.

For a sphere,

\[
\boxed{
S
=
4\pi r^2
}
\]

and

\[
\boxed{
V
=
\frac43\pi r^3.
}
\]

%%%%%%%%%%%%%%%%%%%%%%%%%%%%%%%%%%%%%%%%%%%%%%%%%%%%%%%%%%%%
\subsection{\(\pi\) in Trigonometry}
\label{subsec:pi-trigonometry}
%%%%%%%%%%%%%%%%%%%%%%%%%%%%%%%%%%%%%%%%%%%%%%%%%%%%%%%%%%%%

A full revolution is

\[
\boxed{
2\pi
}
\]

radians.

A half revolution is

\[
\boxed{
\pi,
}
\]

and a quarter revolution is

\[
\boxed{
\frac{\pi}{2}.
}
\]

Thus

\[
\boxed{
180^\circ
=
\pi\text{ radians}.
}
\]

%%%%%%%%%%%%%%%%%%%%%%%%%%%%%%%%%%%%%%%%%%%%%%%%%%%%%%%%%%%%
\subsection{\(\pi\) in Analysis}
\label{subsec:pi-analysis}
%%%%%%%%%%%%%%%%%%%%%%%%%%%%%%%%%%%%%%%%%%%%%%%%%%%%%%%%%%%%

The constant \(\pi\) appears throughout analysis.

For example,

\[
\boxed{
\int_{-\infty}^{\infty}
e^{-x^2}\,dx
=
\sqrt{\pi}.
}
\]

It also appears in Fourier analysis, complex analysis, probability,
differential equations, and geometry.

%%%%%%%%%%%%%%%%%%%%%%%%%%%%%%%%%%%%%%%%%%%%%%%%%%%%%%%%%%%%
\subsection{Irrationality and Transcendence of \(\pi\)}
\label{subsec:pi-irrational-transcendental}
%%%%%%%%%%%%%%%%%%%%%%%%%%%%%%%%%%%%%%%%%%%%%%%%%%%%%%%%%%%%

The number \(\pi\) is irrational:

\[
\boxed{
\pi\notin\mathbb{Q}.
}
\]

It is also transcendental.

Thus \(\pi\) is not a root of any nonzero polynomial with integer
coefficients.

%%%%%%%%%%%%%%%%%%%%%%%%%%%%%%%%%%%%%%%%%%%%%%%%%%%%%%%%%%%%
\subsection{Euler's Number \(e\)}
\label{subsec:euler-number-reference}
%%%%%%%%%%%%%%%%%%%%%%%%%%%%%%%%%%%%%%%%%%%%%%%%%%%%%%%%%%%%

Euler's number is

\[
\boxed{
e
\approx
2.718281828459045.
}
\]

It is the natural base of exponential growth and logarithms.

One defining limit is

\[
\boxed{
e
=
\lim_{n\to\infty}
\left(
1+\frac1n
\right)^n.
}
\]

%%%%%%%%%%%%%%%%%%%%%%%%%%%%%%%%%%%%%%%%%%%%%%%%%%%%%%%%%%%%
\subsection{Series Definition of \(e\)}
\label{subsec:e-series-definition}
%%%%%%%%%%%%%%%%%%%%%%%%%%%%%%%%%%%%%%%%%%%%%%%%%%%%%%%%%%%%

Euler's number also has the series expansion

\[
\boxed{
e
=
\sum_{n=0}^{\infty}
\frac{1}{n!}.
}
\]

Thus

\[
e
=
1
+
1
+
\frac{1}{2!}
+
\frac{1}{3!}
+
\cdots.
\]

%%%%%%%%%%%%%%%%%%%%%%%%%%%%%%%%%%%%%%%%%%%%%%%%%%%%%%%%%%%%
\subsection{The Exponential Function}
\label{subsec:natural-exponential-reference}
%%%%%%%%%%%%%%%%%%%%%%%%%%%%%%%%%%%%%%%%%%%%%%%%%%%%%%%%%%%%

The natural exponential function is

\[
\boxed{
e^x.
}
\]

It is uniquely characterized among exponential functions by

\[
\boxed{
\frac{d}{dx}e^x
=
e^x.
}
\]

Also,

\[
\boxed{
\int e^x\,dx
=
e^x+C.
}
\]

%%%%%%%%%%%%%%%%%%%%%%%%%%%%%%%%%%%%%%%%%%%%%%%%%%%%%%%%%%%%
\subsection{Natural Logarithm}
\label{subsec:natural-logarithm-constant-reference}
%%%%%%%%%%%%%%%%%%%%%%%%%%%%%%%%%%%%%%%%%%%%%%%%%%%%%%%%%%%%

The natural logarithm

\[
\boxed{
\ln x
}
\]

is logarithm base \(e\).

Thus

\[
\boxed{
\ln e=1.
}
\]

It satisfies

\[
\boxed{
\frac{d}{dx}\ln x
=
\frac1x,
\qquad
x>0.
}
\]

%%%%%%%%%%%%%%%%%%%%%%%%%%%%%%%%%%%%%%%%%%%%%%%%%%%%%%%%%%%%
\subsection{Compound Growth and \(e\)}
\label{subsec:e-compound-growth}
%%%%%%%%%%%%%%%%%%%%%%%%%%%%%%%%%%%%%%%%%%%%%%%%%%%%%%%%%%%%

If an amount grows continuously at rate \(r\), then

\[
\boxed{
A(t)
=
A_0e^{rt}.
}
\]

Thus \(e\) appears naturally in:

\[
\boxed{
\text{population growth},
}
\]

\[
\boxed{
\text{radioactive decay},
}
\]

\[
\boxed{
\text{interest},
}
\]

and

\[
\boxed{
\text{differential equations}.
}
\]

%%%%%%%%%%%%%%%%%%%%%%%%%%%%%%%%%%%%%%%%%%%%%%%%%%%%%%%%%%%%
\subsection{Irrationality and Transcendence of \(e\)}
\label{subsec:e-irrational-transcendental}
%%%%%%%%%%%%%%%%%%%%%%%%%%%%%%%%%%%%%%%%%%%%%%%%%%%%%%%%%%%%

Euler's number is irrational and transcendental.

Thus

\[
\boxed{
e\notin\mathbb{Q}.
}
\]

It is also not algebraic over \(\mathbb{Q}\).

%%%%%%%%%%%%%%%%%%%%%%%%%%%%%%%%%%%%%%%%%%%%%%%%%%%%%%%%%%%%
\subsection{Euler's Identity}
\label{subsec:euler-identity-reference}
%%%%%%%%%%%%%%%%%%%%%%%%%%%%%%%%%%%%%%%%%%%%%%%%%%%%%%%%%%%%

One of the most famous formulas in mathematics is

\[
\boxed{
e^{i\pi}+1=0.
}
\]

This identity combines

\[
\boxed{
e,
\quad
i,
\quad
\pi,
\quad
1,
\quad
0.
}
\]

It follows from Euler's formula

\[
\boxed{
e^{i\theta}
=
\cos\theta+i\sin\theta.
}
\]

%%%%%%%%%%%%%%%%%%%%%%%%%%%%%%%%%%%%%%%%%%%%%%%%%%%%%%%%%%%%
\subsection{The Imaginary Unit \(i\)}
\label{subsec:imaginary-unit-reference}
%%%%%%%%%%%%%%%%%%%%%%%%%%%%%%%%%%%%%%%%%%%%%%%%%%%%%%%%%%%%

The imaginary unit is defined by

\[
\boxed{
i^2=-1.
}
\]

Thus

\[
\boxed{
i=\sqrt{-1}
}
\]

under the usual principal-root convention.

It extends the real numbers to the complex numbers.

%%%%%%%%%%%%%%%%%%%%%%%%%%%%%%%%%%%%%%%%%%%%%%%%%%%%%%%%%%%%
\subsection{Powers of \(i\)}
\label{subsec:powers-of-i-reference}
%%%%%%%%%%%%%%%%%%%%%%%%%%%%%%%%%%%%%%%%%%%%%%%%%%%%%%%%%%%%

The powers of \(i\) repeat with period \(4\):

\[
\boxed{
i^0=1,
}
\]

\[
\boxed{
i^1=i,
}
\]

\[
\boxed{
i^2=-1,
}
\]

\[
\boxed{
i^3=-i,
}
\]

\[
\boxed{
i^4=1.
}
\]

Therefore,

\[
\boxed{
i^{n+4}=i^n.
}
\]

%%%%%%%%%%%%%%%%%%%%%%%%%%%%%%%%%%%%%%%%%%%%%%%%%%%%%%%%%%%%
\subsection{The Golden Ratio}
\label{subsec:golden-ratio-reference}
%%%%%%%%%%%%%%%%%%%%%%%%%%%%%%%%%%%%%%%%%%%%%%%%%%%%%%%%%%%%

The golden ratio is

\[
\boxed{
\phi
=
\frac{1+\sqrt5}{2}.
}
\]

Its decimal expansion begins

\[
\boxed{
\phi
\approx
1.618033988749895.
}
\]

It is algebraic because it satisfies

\[
\boxed{
\phi^2-\phi-1=0.
}
\]

%%%%%%%%%%%%%%%%%%%%%%%%%%%%%%%%%%%%%%%%%%%%%%%%%%%%%%%%%%%%
\subsection{Basic Golden Ratio Identities}
\label{subsec:golden-ratio-identities}
%%%%%%%%%%%%%%%%%%%%%%%%%%%%%%%%%%%%%%%%%%%%%%%%%%%%%%%%%%%%

From

\[
\phi^2=\phi+1,
\]

we obtain

\[
\boxed{
\phi-1
=
\frac1\phi.
}
\]

Also,

\[
\boxed{
\phi^2
=
\phi+1.
}
\]

The conjugate root of

\[
x^2-x-1=0
\]

is

\[
\boxed{
\psi
=
\frac{1-\sqrt5}{2}.
}
\]

%%%%%%%%%%%%%%%%%%%%%%%%%%%%%%%%%%%%%%%%%%%%%%%%%%%%%%%%%%%%
\subsection{Golden Ratio and Fibonacci Numbers}
\label{subsec:golden-ratio-fibonacci}
%%%%%%%%%%%%%%%%%%%%%%%%%%%%%%%%%%%%%%%%%%%%%%%%%%%%%%%%%%%%

Let Fibonacci numbers satisfy

\[
F_{n+1}
=
F_n+F_{n-1}.
\]

Then

\[
\boxed{
\frac{F_{n+1}}{F_n}
\to
\phi
}
\]

as

\[
n\to\infty.
\]

Binet's formula is

\[
\boxed{
F_n
=
\frac{
\phi^n-\psi^n
}{
\sqrt5
},
}
\]

where

\[
\psi
=
\frac{1-\sqrt5}{2}.
\]

%%%%%%%%%%%%%%%%%%%%%%%%%%%%%%%%%%%%%%%%%%%%%%%%%%%%%%%%%%%%
\subsection{Euler--Mascheroni Constant}
\label{subsec:euler-mascheroni-reference}
%%%%%%%%%%%%%%%%%%%%%%%%%%%%%%%%%%%%%%%%%%%%%%%%%%%%%%%%%%%%

The Euler--Mascheroni constant is

\[
\boxed{
\gamma
\approx
0.5772156649015329.
}
\]

It is defined by

\[
\boxed{
\gamma
=
\lim_{n\to\infty}
\left(
\sum_{k=1}^{n}\frac1k
-
\ln n
\right).
}
\]

It measures the asymptotic difference between the harmonic series and the
natural logarithm.

%%%%%%%%%%%%%%%%%%%%%%%%%%%%%%%%%%%%%%%%%%%%%%%%%%%%%%%%%%%%
\subsection{Harmonic Numbers}
\label{subsec:harmonic-numbers-constant-reference}
%%%%%%%%%%%%%%%%%%%%%%%%%%%%%%%%%%%%%%%%%%%%%%%%%%%%%%%%%%%%

The \(n\)-th harmonic number is

\[
\boxed{
H_n
=
\sum_{k=1}^{n}
\frac1k.
}
\]

As \(n\to\infty\),

\[
\boxed{
H_n
=
\ln n
+
\gamma
+
O\left(\frac1n\right).
}
\]

A more refined expansion is

\[
\boxed{
H_n
=
\ln n
+
\gamma
+
\frac{1}{2n}
-
\frac{1}{12n^2}
+
O(n^{-4}).
}
\]

%%%%%%%%%%%%%%%%%%%%%%%%%%%%%%%%%%%%%%%%%%%%%%%%%%%%%%%%%%%%
\subsection{The Square Root of Two}
\label{subsec:sqrt2-constant-reference}
%%%%%%%%%%%%%%%%%%%%%%%%%%%%%%%%%%%%%%%%%%%%%%%%%%%%%%%%%%%%

The constant

\[
\boxed{
\sqrt2
\approx
1.414213562373095
}
\]

is the diagonal length of a unit square.

By the Pythagorean theorem,

\[
\boxed{
d
=
\sqrt{1^2+1^2}
=
\sqrt2.
}
\]

It is irrational but algebraic.

%%%%%%%%%%%%%%%%%%%%%%%%%%%%%%%%%%%%%%%%%%%%%%%%%%%%%%%%%%%%
\subsection{The Square Root of Three}
\label{subsec:sqrt3-constant-reference}
%%%%%%%%%%%%%%%%%%%%%%%%%%%%%%%%%%%%%%%%%%%%%%%%%%%%%%%%%%%%

The constant

\[
\boxed{
\sqrt3
\approx
1.732050807568877
}
\]

appears naturally in equilateral triangles and trigonometry.

For an equilateral triangle of side \(2\), the altitude is

\[
\boxed{
\sqrt3.
}
\]

It is irrational and algebraic.

%%%%%%%%%%%%%%%%%%%%%%%%%%%%%%%%%%%%%%%%%%%%%%%%%%%%%%%%%%%%
\subsection{The Natural Logarithm of Two}
\label{subsec:ln2-constant-reference}
%%%%%%%%%%%%%%%%%%%%%%%%%%%%%%%%%%%%%%%%%%%%%%%%%%%%%%%%%%%%

The constant

\[
\boxed{
\ln2
\approx
0.693147180559945.
}
\]

It has the integral representation

\[
\boxed{
\ln2
=
\int_1^2
\frac{1}{x}\,dx.
}
\]

It also appears in alternating harmonic-series calculations.

%%%%%%%%%%%%%%%%%%%%%%%%%%%%%%%%%%%%%%%%%%%%%%%%%%%%%%%%%%%%
\subsection{Alternating Harmonic Series}
\label{subsec:alternating-harmonic-constant}
%%%%%%%%%%%%%%%%%%%%%%%%%%%%%%%%%%%%%%%%%%%%%%%%%%%%%%%%%%%%

The alternating harmonic series satisfies

\[
\boxed{
1
-
\frac12
+
\frac13
-
\frac14
+
\cdots
=
\ln2.
}
\]

Equivalently,

\[
\boxed{
\ln2
=
\sum_{n=1}^{\infty}
\frac{(-1)^{n+1}}{n}.
}
\]

%%%%%%%%%%%%%%%%%%%%%%%%%%%%%%%%%%%%%%%%%%%%%%%%%%%%%%%%%%%%
\subsection{The Basel Constant}
\label{subsec:basel-constant-reference}
%%%%%%%%%%%%%%%%%%%%%%%%%%%%%%%%%%%%%%%%%%%%%%%%%%%%%%%%%%%%

The sum

\[
\sum_{n=1}^{\infty}
\frac{1}{n^2}
\]

equals

\[
\boxed{
\frac{\pi^2}{6}.
}
\]

Thus

\[
\boxed{
\zeta(2)
=
\frac{\pi^2}{6}
\approx
1.644934066848226.
}
\]

%%%%%%%%%%%%%%%%%%%%%%%%%%%%%%%%%%%%%%%%%%%%%%%%%%%%%%%%%%%%
\subsection{Riemann Zeta Values}
\label{subsec:zeta-values-reference}
%%%%%%%%%%%%%%%%%%%%%%%%%%%%%%%%%%%%%%%%%%%%%%%%%%%%%%%%%%%%

The Riemann zeta function is

\[
\boxed{
\zeta(s)
=
\sum_{n=1}^{\infty}
\frac{1}{n^s}
}
\]

for

\[
\operatorname{Re}(s)>1.
\]

Important values include

\[
\boxed{
\zeta(2)
=
\frac{\pi^2}{6},
}
\]

and

\[
\boxed{
\zeta(4)
=
\frac{\pi^4}{90}.
}
\]

More generally,

\[
\zeta(2n)
\]

can be expressed using powers of \(\pi\) and Bernoulli numbers.

%%%%%%%%%%%%%%%%%%%%%%%%%%%%%%%%%%%%%%%%%%%%%%%%%%%%%%%%%%%%
\subsection{Apery's Constant}
\label{subsec:apery-constant-reference}
%%%%%%%%%%%%%%%%%%%%%%%%%%%%%%%%%%%%%%%%%%%%%%%%%%%%%%%%%%%%

Apery's constant is

\[
\boxed{
\zeta(3)
=
\sum_{n=1}^{\infty}
\frac{1}{n^3}.
}
\]

Numerically,

\[
\boxed{
\zeta(3)
\approx
1.202056903159594.
}
\]

It is known to be irrational.

%%%%%%%%%%%%%%%%%%%%%%%%%%%%%%%%%%%%%%%%%%%%%%%%%%%%%%%%%%%%
\subsection{Catalan's Constant}
\label{subsec:catalan-constant-reference}
%%%%%%%%%%%%%%%%%%%%%%%%%%%%%%%%%%%%%%%%%%%%%%%%%%%%%%%%%%%%

Catalan's constant is

\[
\boxed{
G
=
\sum_{n=0}^{\infty}
\frac{(-1)^n}{(2n+1)^2}.
}
\]

Its approximate value is

\[
\boxed{
G
\approx
0.915965594177219.
}
\]

It appears in combinatorics, analysis, geometry, and special-function
integrals.

%%%%%%%%%%%%%%%%%%%%%%%%%%%%%%%%%%%%%%%%%%%%%%%%%%%%%%%%%%%%
\subsection{Catalan Constant Integral}
\label{subsec:catalan-integral-reference}
%%%%%%%%%%%%%%%%%%%%%%%%%%%%%%%%%%%%%%%%%%%%%%%%%%%%%%%%%%%%

One representation is

\[
\boxed{
G
=
-\int_0^{\pi/4}
\ln(\tan x)\,dx.
}
\]

This illustrates how constants can arise from definite integrals even when
they do not have simple elementary closed forms.

%%%%%%%%%%%%%%%%%%%%%%%%%%%%%%%%%%%%%%%%%%%%%%%%%%%%%%%%%%%%
\subsection{Feigenbaum Constant}
\label{subsec:feigenbaum-constant-reference}
%%%%%%%%%%%%%%%%%%%%%%%%%%%%%%%%%%%%%%%%%%%%%%%%%%%%%%%%%%%%

A major constant in nonlinear dynamics is the Feigenbaum constant

\[
\boxed{
\delta
\approx
4.669201609102990.
}
\]

It describes the asymptotic ratio of parameter intervals in period-doubling
routes to chaos for a broad class of one-dimensional dynamical systems.

%%%%%%%%%%%%%%%%%%%%%%%%%%%%%%%%%%%%%%%%%%%%%%%%%%%%%%%%%%%%
\subsection{Feigenbaum Scaling}
\label{subsec:feigenbaum-scaling}
%%%%%%%%%%%%%%%%%%%%%%%%%%%%%%%%%%%%%%%%%%%%%%%%%%%%%%%%%%%%

If

\[
r_1,r_2,r_3,\ldots
\]

are successive parameter values at which period doubling occurs, then under
appropriate universality conditions,

\[
\boxed{
\frac{
r_{n}-r_{n-1}
}{
r_{n+1}-r_n
}
\to
\delta.
}
\]

This is an example of a universal constant arising from dynamical behavior.

%%%%%%%%%%%%%%%%%%%%%%%%%%%%%%%%%%%%%%%%%%%%%%%%%%%%%%%%%%%%
\subsection{Khinchin's Constant}
\label{subsec:khinchin-constant-reference}
%%%%%%%%%%%%%%%%%%%%%%%%%%%%%%%%%%%%%%%%%%%%%%%%%%%%%%%%%%%%

Khinchin's constant is approximately

\[
\boxed{
K_0
\approx
2.685452001065306.
}
\]

It appears in continued fractions.

For almost every real number under the relevant measure-theoretic sense, the
geometric mean of continued-fraction coefficients converges to \(K_0\).

%%%%%%%%%%%%%%%%%%%%%%%%%%%%%%%%%%%%%%%%%%%%%%%%%%%%%%%%%%%%
\subsection{Glaisher--Kinkelin Constant}
\label{subsec:glaisher-kinkelin-reference}
%%%%%%%%%%%%%%%%%%%%%%%%%%%%%%%%%%%%%%%%%%%%%%%%%%%%%%%%%%%%

The Glaisher--Kinkelin constant is commonly denoted

\[
\boxed{
A
}
\]

and satisfies approximately

\[
\boxed{
A
\approx
1.282427129100623.
}
\]

It arises in products involving powers of integers, special functions, and
asymptotic formulas.

%%%%%%%%%%%%%%%%%%%%%%%%%%%%%%%%%%%%%%%%%%%%%%%%%%%%%%%%%%%%
\subsection{Omega Constant}
\label{subsec:omega-constant-reference}
%%%%%%%%%%%%%%%%%%%%%%%%%%%%%%%%%%%%%%%%%%%%%%%%%%%%%%%%%%%%

The omega constant is the unique positive real number satisfying

\[
\boxed{
\Omega e^\Omega=1.
}
\]

Equivalently,

\[
\boxed{
\Omega
=
e^{-\Omega}.
}
\]

Its approximate value is

\[
\boxed{
\Omega
\approx
0.567143290409784.
}
\]

It is related to the Lambert \(W\) function:

\[
\boxed{
\Omega
=
W(1).
}
\]

%%%%%%%%%%%%%%%%%%%%%%%%%%%%%%%%%%%%%%%%%%%%%%%%%%%%%%%%%%%%
\subsection{Plastic Constant}
\label{subsec:plastic-constant-reference}
%%%%%%%%%%%%%%%%%%%%%%%%%%%%%%%%%%%%%%%%%%%%%%%%%%%%%%%%%%%%

The plastic constant is the unique real root greater than \(1\) of

\[
\boxed{
x^3=x+1.
}
\]

It is approximately

\[
\boxed{
\rho_{\mathrm{plastic}}
\approx
1.324717957244746.
}
\]

It appears in algebraic recurrences and certain geometric constructions.

%%%%%%%%%%%%%%%%%%%%%%%%%%%%%%%%%%%%%%%%%%%%%%%%%%%%%%%%%%%%
\subsection{Euler's Constant Versus Euler's Number}
\label{subsec:euler-constant-versus-number}
%%%%%%%%%%%%%%%%%%%%%%%%%%%%%%%%%%%%%%%%%%%%%%%%%%%%%%%%%%%%

Do not confuse

\[
\boxed{
e
\approx
2.71828
}
\]

with

\[
\boxed{
\gamma
\approx
0.57722.
}
\]

Euler's number \(e\) is the natural exponential base.

The Euler--Mascheroni constant \(\gamma\) is the asymptotic difference between
harmonic numbers and logarithms.

%%%%%%%%%%%%%%%%%%%%%%%%%%%%%%%%%%%%%%%%%%%%%%%%%%%%%%%%%%%%
\subsection{Golden Ratio Versus Euler Totient}
\label{subsec:phi-ambiguity-constants}
%%%%%%%%%%%%%%%%%%%%%%%%%%%%%%%%%%%%%%%%%%%%%%%%%%%%%%%%%%%%

The symbol

\[
\phi
\]

may denote the golden ratio:

\[
\boxed{
\phi
=
\frac{1+\sqrt5}{2},
}
\]

while in number theory,

\[
\boxed{
\phi(n)
}
\]

or

\[
\boxed{
\varphi(n)
}
\]

may denote Euler's totient function.

Context distinguishes the meanings.

%%%%%%%%%%%%%%%%%%%%%%%%%%%%%%%%%%%%%%%%%%%%%%%%%%%%%%%%%%%%
\subsection{Constants in Geometry}
\label{subsec:constants-in-geometry}
%%%%%%%%%%%%%%%%%%%%%%%%%%%%%%%%%%%%%%%%%%%%%%%%%%%%%%%%%%%%

Important geometric constants include

\[
\boxed{
\pi,
}
\]

\[
\boxed{
\sqrt2,
}
\]

\[
\boxed{
\sqrt3,
}
\]

and

\[
\boxed{
\phi.
}
\]

They arise from:

\[
\boxed{
\text{circles},
}
\]

\[
\boxed{
\text{diagonals},
}
\]

\[
\boxed{
\text{regular polygons},
}
\]

and

\[
\boxed{
\text{similarity ratios}.
}
\]

%%%%%%%%%%%%%%%%%%%%%%%%%%%%%%%%%%%%%%%%%%%%%%%%%%%%%%%%%%%%
\subsection{Constants in Calculus}
\label{subsec:constants-in-calculus}
%%%%%%%%%%%%%%%%%%%%%%%%%%%%%%%%%%%%%%%%%%%%%%%%%%%%%%%%%%%%

Calculus frequently uses

\[
\boxed{
e
}
\]

and

\[
\boxed{
\pi.
}
\]

Examples include

\[
\boxed{
\frac{d}{dx}e^x=e^x
}
\]

and

\[
\boxed{
\int_{-\infty}^{\infty}
e^{-x^2}\,dx
=
\sqrt{\pi}.
}
\]

%%%%%%%%%%%%%%%%%%%%%%%%%%%%%%%%%%%%%%%%%%%%%%%%%%%%%%%%%%%%
\subsection{Constants in Complex Analysis}
\label{subsec:constants-in-complex-analysis}
%%%%%%%%%%%%%%%%%%%%%%%%%%%%%%%%%%%%%%%%%%%%%%%%%%%%%%%%%%%%

The constants

\[
\boxed{
e,
\quad
i,
\quad
\pi
}
\]

are connected through

\[
\boxed{
e^{i\theta}
=
\cos\theta+i\sin\theta.
}
\]

At

\[
\theta=\pi,
\]

this gives

\[
\boxed{
e^{i\pi}=-1.
}
\]

%%%%%%%%%%%%%%%%%%%%%%%%%%%%%%%%%%%%%%%%%%%%%%%%%%%%%%%%%%%%
\subsection{Constants in Probability}
\label{subsec:constants-in-probability}
%%%%%%%%%%%%%%%%%%%%%%%%%%%%%%%%%%%%%%%%%%%%%%%%%%%%%%%%%%%%

The normal density contains both \(e\) and \(\pi\):

\[
\boxed{
f(x)
=
\frac{1}{\sigma\sqrt{2\pi}}
\exp
\left(
-\frac{(x-\mu)^2}{2\sigma^2}
\right).
}
\]

Thus fundamental constants of analysis arise directly in probability.

%%%%%%%%%%%%%%%%%%%%%%%%%%%%%%%%%%%%%%%%%%%%%%%%%%%%%%%%%%%%
\subsection{The Standard Normal Constant}
\label{subsec:standard-normal-constant}
%%%%%%%%%%%%%%%%%%%%%%%%%%%%%%%%%%%%%%%%%%%%%%%%%%%%%%%%%%%%

For

\[
Z\sim N(0,1),
\]

the density is

\[
\boxed{
\phi_Z(z)
=
\frac{1}{\sqrt{2\pi}}
e^{-z^2/2}.
}
\]

The normalization constant

\[
\boxed{
\frac{1}{\sqrt{2\pi}}
}
\]

ensures that the density integrates to \(1\).

%%%%%%%%%%%%%%%%%%%%%%%%%%%%%%%%%%%%%%%%%%%%%%%%%%%%%%%%%%%%
\subsection{Constants in Fourier Analysis}
\label{subsec:constants-in-fourier-analysis}
%%%%%%%%%%%%%%%%%%%%%%%%%%%%%%%%%%%%%%%%%%%%%%%%%%%%%%%%%%%%

Fourier-transform conventions often involve factors such as

\[
\boxed{
2\pi
}
\]

or

\[
\boxed{
\frac{1}{\sqrt{2\pi}}.
}
\]

For example, one convention is

\[
\boxed{
\hat f(\omega)
=
\int_{-\infty}^{\infty}
f(x)e^{-i\omega x}\,dx.
}
\]

Other conventions redistribute \(2\pi\) factors between the forward and
inverse transforms.

%%%%%%%%%%%%%%%%%%%%%%%%%%%%%%%%%%%%%%%%%%%%%%%%%%%%%%%%%%%%
\subsection{Constants in Differential Equations}
\label{subsec:constants-in-differential-equations}
%%%%%%%%%%%%%%%%%%%%%%%%%%%%%%%%%%%%%%%%%%%%%%%%%%%%%%%%%%%%

The constants \(e\) and \(\pi\) appear naturally in differential equations.

For

\[
y'=ky,
\]

the solution is

\[
\boxed{
y(t)
=
Ce^{kt}.
}
\]

For the harmonic oscillator

\[
y''+\omega^2y=0,
\]

period is

\[
\boxed{
T
=
\frac{2\pi}{\omega}.
}
\]

%%%%%%%%%%%%%%%%%%%%%%%%%%%%%%%%%%%%%%%%%%%%%%%%%%%%%%%%%%%%
\subsection{Constants in Number Theory}
\label{subsec:constants-in-number-theory}
%%%%%%%%%%%%%%%%%%%%%%%%%%%%%%%%%%%%%%%%%%%%%%%%%%%%%%%%%%%%

Number theory contains many important constants, including

\[
\boxed{
\pi,
}
\]

\[
\boxed{
e,
}
\]

\[
\boxed{
\gamma,
}
\]

and

\[
\boxed{
\zeta(3).
}
\]

The prime number theorem contains the asymptotic relation

\[
\boxed{
\pi(x)
\sim
\frac{x}{\ln x},
}
\]

where here \(\pi(x)\) denotes the prime-counting function rather than the
circle constant.

%%%%%%%%%%%%%%%%%%%%%%%%%%%%%%%%%%%%%%%%%%%%%%%%%%%%%%%%%%%%
\subsection{Constants in Combinatorics}
\label{subsec:constants-in-combinatorics}
%%%%%%%%%%%%%%%%%%%%%%%%%%%%%%%%%%%%%%%%%%%%%%%%%%%%%%%%%%%%

The number \(e\) appears in counting problems.

For example, the number of derangements \(D_n\) satisfies

\[
\boxed{
D_n
=
n!
\sum_{k=0}^{n}
\frac{(-1)^k}{k!}.
}
\]

Asymptotically,

\[
\boxed{
D_n
\sim
\frac{n!}{e}.
}
\]

%%%%%%%%%%%%%%%%%%%%%%%%%%%%%%%%%%%%%%%%%%%%%%%%%%%%%%%%%%%%
\subsection{Constants in Asymptotic Analysis}
\label{subsec:constants-in-asymptotic-analysis}
%%%%%%%%%%%%%%%%%%%%%%%%%%%%%%%%%%%%%%%%%%%%%%%%%%%%%%%%%%%%

Stirling's approximation connects \(e\), \(\pi\), and factorials:

\[
\boxed{
n!
\sim
\sqrt{2\pi n}
\left(
\frac{n}{e}
\right)^n.
}
\]

This is one of the most important asymptotic formulas in mathematics.

%%%%%%%%%%%%%%%%%%%%%%%%%%%%%%%%%%%%%%%%%%%%%%%%%%%%%%%%%%%%
\subsection{Worked Example: Stirling Approximation}
\label{subsec:worked-stirling-approximation}
%%%%%%%%%%%%%%%%%%%%%%%%%%%%%%%%%%%%%%%%%%%%%%%%%%%%%%%%%%%%

\begin{workedexamplebox}{Approximate \(10!\)}

Using Stirling's formula,

\[
10!
\approx
\sqrt{20\pi}
\left(
\frac{10}{e}
\right)^{10}.
\]

Numerically,

\[
\sqrt{20\pi}
\left(
\frac{10}{e}
\right)^{10}
\approx
3.5987\times10^6.
\]

The exact value is

\[
10!
=
3.6288\times10^6.
\]

\begin{answerbox}
\[
\boxed{
10!
\approx
3.60\times10^6
}
\]

using the leading Stirling approximation.
\end{answerbox}

\end{workedexamplebox}

%%%%%%%%%%%%%%%%%%%%%%%%%%%%%%%%%%%%%%%%%%%%%%%%%%%%%%%%%%%%
\subsection{Constants in Dynamical Systems}
\label{subsec:constants-in-dynamical-systems}
%%%%%%%%%%%%%%%%%%%%%%%%%%%%%%%%%%%%%%%%%%%%%%%%%%%%%%%%%%%%

Dynamical systems contain universal constants such as the Feigenbaum
constants.

These arise from repeated scaling patterns near transitions to chaos.

Unlike geometric constants such as \(\pi\), they emerge from asymptotic
behavior of families of nonlinear maps.

%%%%%%%%%%%%%%%%%%%%%%%%%%%%%%%%%%%%%%%%%%%%%%%%%%%%%%%%%%%%
\subsection{Constants in Optimization}
\label{subsec:constants-in-optimization}
%%%%%%%%%%%%%%%%%%%%%%%%%%%%%%%%%%%%%%%%%%%%%%%%%%%%%%%%%%%%

Optimization often contains fixed numerical constants inside convergence
bounds.

For example, for a strongly convex quadratic problem, rates may depend on a
condition number

\[
\boxed{
\kappa
=
\frac{L}{\mu}.
}
\]

Here \(L\) and \(\mu\) are parameters rather than universal constants.

This distinction is important:

\[
\boxed{
\text{constant within a problem}
\neq
\text{universal mathematical constant}.
}
\]

%%%%%%%%%%%%%%%%%%%%%%%%%%%%%%%%%%%%%%%%%%%%%%%%%%%%%%%%%%%%
\subsection{Universal Constants Versus Parameters}
\label{subsec:universal-constants-versus-parameters}
%%%%%%%%%%%%%%%%%%%%%%%%%%%%%%%%%%%%%%%%%%%%%%%%%%%%%%%%%%%%

A universal mathematical constant such as

\[
\pi
\]

has the same value in every mathematical problem.

A parameter such as

\[
\lambda
\]

may be fixed within a model but can change from one problem to another.

Thus the word constant can mean either:

\[
\boxed{
\text{universally fixed numerical value}
}
\]

or

\[
\boxed{
\text{quantity held fixed within a calculation}.
}
\]

%%%%%%%%%%%%%%%%%%%%%%%%%%%%%%%%%%%%%%%%%%%%%%%%%%%%%%%%%%%%
\subsection{Constants of Integration}
\label{subsec:constants-of-integration-reference}
%%%%%%%%%%%%%%%%%%%%%%%%%%%%%%%%%%%%%%%%%%%%%%%%%%%%%%%%%%%%

When taking an indefinite integral,

\[
\boxed{
\int f(x)\,dx
=
F(x)+C.
}
\]

The symbol \(C\) is called the constant of integration.

It is not one universal number.

Rather, it represents an arbitrary real constant because

\[
\frac{d}{dx}C=0.
\]

%%%%%%%%%%%%%%%%%%%%%%%%%%%%%%%%%%%%%%%%%%%%%%%%%%%%%%%%%%%%
\subsection{Worked Example: Constant of Integration}
\label{subsec:worked-constant-integration}
%%%%%%%%%%%%%%%%%%%%%%%%%%%%%%%%%%%%%%%%%%%%%%%%%%%%%%%%%%%%

\begin{workedexamplebox}{Antiderivative Family}

Compute

\[
\int 2x\,dx.
\]

Since

\[
\frac{d}{dx}x^2=2x,
\]

we have

\[
\begin{answerbox}
\[
\boxed{
\int2x\,dx
=
x^2+C.
}
\]
\end{answerbox}

Every value of \(C\) gives a different antiderivative with the same
derivative.

\end{workedexamplebox}

%%%%%%%%%%%%%%%%%%%%%%%%%%%%%%%%%%%%%%%%%%%%%%%%%%%%%%%%%%%%
\subsection{Physical Constants Versus Mathematical Constants}
\label{subsec:physical-versus-mathematical-constants}
%%%%%%%%%%%%%%%%%%%%%%%%%%%%%%%%%%%%%%%%%%%%%%%%%%%%%%%%%%%%

Mathematical constants are determined by mathematical definitions.

Physical constants are experimentally measured quantities describing the
physical universe.

Examples of mathematical constants include

\[
\boxed{
\pi,
\quad
e,
\quad
\gamma.
}
\]

Examples of physical constants include quantities such as the speed of light
or gravitational constant.

Physical constants carry units or depend on unit conventions.

Pure mathematical constants are generally dimensionless.

%%%%%%%%%%%%%%%%%%%%%%%%%%%%%%%%%%%%%%%%%%%%%%%%%%%%%%%%%%%%
\subsection{Algebraic Constants}
\label{subsec:algebraic-constants-reference}
%%%%%%%%%%%%%%%%%%%%%%%%%%%%%%%%%%%%%%%%%%%%%%%%%%%%%%%%%%%%

A constant is algebraic if it satisfies a polynomial equation with integer
coefficients.

Examples include

\[
\boxed{
\sqrt2
}
\]

because

\[
x^2-2=0,
\]

and

\[
\boxed{
\phi
}
\]

because

\[
x^2-x-1=0.
\]

%%%%%%%%%%%%%%%%%%%%%%%%%%%%%%%%%%%%%%%%%%%%%%%%%%%%%%%%%%%%
\subsection{Transcendental Constants}
\label{subsec:transcendental-constants-reference}
%%%%%%%%%%%%%%%%%%%%%%%%%%%%%%%%%%%%%%%%%%%%%%%%%%%%%%%%%%%%

A transcendental constant is not algebraic.

Important examples include

\[
\boxed{
\pi
}
\]

and

\[
\boxed{
e.
}
\]

Thus transcendental numbers cannot be roots of nonzero integer-coefficient
polynomials.

%%%%%%%%%%%%%%%%%%%%%%%%%%%%%%%%%%%%%%%%%%%%%%%%%%%%%%%%%%%%
\subsection{Known and Unknown Arithmetic Status}
\label{subsec:constant-arithmetic-status}
%%%%%%%%%%%%%%%%%%%%%%%%%%%%%%%%%%%%%%%%%%%%%%%%%%%%%%%%%%%%

Not every famous mathematical constant has a completely known arithmetic
classification.

For example, some constants are known to be irrational but not known to be
transcendental, while for others even irrationality remains unresolved.

Therefore one should not assume that every named constant has a known
algebraic classification.

%%%%%%%%%%%%%%%%%%%%%%%%%%%%%%%%%%%%%%%%%%%%%%%%%%%%%%%%%%%%
\subsection{Approximating Constants Numerically}
\label{subsec:approximating-constants}
%%%%%%%%%%%%%%%%%%%%%%%%%%%%%%%%%%%%%%%%%%%%%%%%%%%%%%%%%%%%

Constants may be approximated using:

\[
\boxed{
\text{series},
}
\]

\[
\boxed{
\text{products},
}
\]

\[
\boxed{
\text{integrals},
}
\]

\[
\boxed{
\text{iterations},
}
\]

and

\[
\boxed{
\text{numerical algorithms}.
}
\]

For example,

\[
e
=
\sum_{n=0}^{\infty}
\frac{1}{n!}
\]

provides a rapidly convergent approximation.

%%%%%%%%%%%%%%%%%%%%%%%%%%%%%%%%%%%%%%%%%%%%%%%%%%%%%%%%%%%%
\subsection{A Series for \(\pi\)}
\label{subsec:pi-series-reference}
%%%%%%%%%%%%%%%%%%%%%%%%%%%%%%%%%%%%%%%%%%%%%%%%%%%%%%%%%%%%

The Gregory--Leibniz series is

\[
\boxed{
\frac{\pi}{4}
=
1
-
\frac13
+
\frac15
-
\frac17
+
\cdots.
}
\]

Equivalently,

\[
\boxed{
\pi
=
4
\sum_{n=0}^{\infty}
\frac{(-1)^n}{2n+1}.
}
\]

This series converges slowly but is mathematically important.

%%%%%%%%%%%%%%%%%%%%%%%%%%%%%%%%%%%%%%%%%%%%%%%%%%%%%%%%%%%%
\subsection{An Integral for \(\pi\)}
\label{subsec:pi-integral-reference}
%%%%%%%%%%%%%%%%%%%%%%%%%%%%%%%%%%%%%%%%%%%%%%%%%%%%%%%%%%%%

Since

\[
\int\frac{1}{1+x^2}\,dx
=
\arctan x,
\]

we obtain

\[
\boxed{
\int_0^1
\frac{1}{1+x^2}\,dx
=
\frac{\pi}{4}.
}
\]

Therefore,

\[
\boxed{
\pi
=
4
\int_0^1
\frac{1}{1+x^2}\,dx.
}
\]

%%%%%%%%%%%%%%%%%%%%%%%%%%%%%%%%%%%%%%%%%%%%%%%%%%%%%%%%%%%%
\subsection{An Integral for \(e\)}
\label{subsec:e-integral-reference}
%%%%%%%%%%%%%%%%%%%%%%%%%%%%%%%%%%%%%%%%%%%%%%%%%%%%%%%%%%%%

Since

\[
\ln e=1,
\]

and

\[
\ln x
=
\int_1^x
\frac{1}{t}\,dt,
\]

the number \(e\) is the unique positive number satisfying

\[
\boxed{
\int_1^e
\frac{1}{t}\,dt
=
1.
}
\]

%%%%%%%%%%%%%%%%%%%%%%%%%%%%%%%%%%%%%%%%%%%%%%%%%%%%%%%%%%%%
\subsection{Constants From Infinite Products}
\label{subsec:constants-infinite-products}
%%%%%%%%%%%%%%%%%%%%%%%%%%%%%%%%%%%%%%%%%%%%%%%%%%%%%%%%%%%%

Important constants also arise from infinite products.

For example, Wallis' product states

\[
\boxed{
\frac{\pi}{2}
=
\prod_{n=1}^{\infty}
\frac{(2n)^2}
{(2n-1)(2n+1)}.
}
\]

Thus

\[
\boxed{
\frac{\pi}{2}
=
\frac21
\cdot
\frac23
\cdot
\frac43
\cdot
\frac45
\cdot
\frac65
\cdot
\frac67
\cdots.
}
\]

%%%%%%%%%%%%%%%%%%%%%%%%%%%%%%%%%%%%%%%%%%%%%%%%%%%%%%%%%%%%
\subsection{Constants From Limits}
\label{subsec:constants-from-limits}
%%%%%%%%%%%%%%%%%%%%%%%%%%%%%%%%%%%%%%%%%%%%%%%%%%%%%%%%%%%%

Many constants can be defined by limits.

Examples include

\[
\boxed{
e
=
\lim_{n\to\infty}
\left(
1+\frac1n
\right)^n
}
\]

and

\[
\boxed{
\gamma
=
\lim_{n\to\infty}
\left(
H_n-\ln n
\right).
}
\]

Limits therefore provide exact definitions of quantities whose decimal
expansions never terminate.

%%%%%%%%%%%%%%%%%%%%%%%%%%%%%%%%%%%%%%%%%%%%%%%%%%%%%%%%%%%%
\subsection{Constants From Optimization}
\label{subsec:constants-from-optimization}
%%%%%%%%%%%%%%%%%%%%%%%%%%%%%%%%%%%%%%%%%%%%%%%%%%%%%%%%%%%%

Some constants can arise as extremal values.

For example, geometric constants can result from minimizing or maximizing
ratios under geometric constraints.

Constants may therefore emerge naturally from

\[
\boxed{
\text{optimal shapes},
}
\]

\[
\boxed{
\text{best approximation},
}
\]

or

\[
\boxed{
\text{extremal inequalities}.
}
\]

%%%%%%%%%%%%%%%%%%%%%%%%%%%%%%%%%%%%%%%%%%%%%%%%%%%%%%%%%%%%
\subsection{Constants and Dimensionless Quantities}
\label{subsec:constants-dimensionless}
%%%%%%%%%%%%%%%%%%%%%%%%%%%%%%%%%%%%%%%%%%%%%%%%%%%%%%%%%%%%

Pure mathematical constants such as

\[
\pi
\]

and

\[
e
\]

are dimensionless.

For example,

\[
\frac{
\text{circumference}
}{
\text{diameter}
}
\]

is a ratio of two lengths, so the units cancel.

This is one reason universal constants often arise from ratios.

%%%%%%%%%%%%%%%%%%%%%%%%%%%%%%%%%%%%%%%%%%%%%%%%%%%%%%%%%%%%
\subsection{Rounding Constants}
\label{subsec:rounding-constants}
%%%%%%%%%%%%%%%%%%%%%%%%%%%%%%%%%%%%%%%%%%%%%%%%%%%%%%%%%%%%

When numerical approximation is required, the number of retained digits
should match the precision of the problem.

For example,

\[
\boxed{
\pi
\approx
3.14
}
\]

may be sufficient for a rough estimate.

For higher precision,

\[
\boxed{
\pi
\approx
3.14159265.
}
\]

Unnecessary decimal digits can create a false impression of precision.

%%%%%%%%%%%%%%%%%%%%%%%%%%%%%%%%%%%%%%%%%%%%%%%%%%%%%%%%%%%%
\subsection{Exact Versus Approximate Calculation}
\label{subsec:exact-approximate-calculation}
%%%%%%%%%%%%%%%%%%%%%%%%%%%%%%%%%%%%%%%%%%%%%%%%%%%%%%%%%%%%

Suppose a circle has radius

\[
r=3.
\]

Its exact area is

\[
\boxed{
A=9\pi.
}
\]

A decimal approximation is

\[
\boxed{
A
\approx
28.2743.
}
\]

The exact form is normally preferable until a decimal answer is required.

%%%%%%%%%%%%%%%%%%%%%%%%%%%%%%%%%%%%%%%%%%%%%%%%%%%%%%%%%%%%
\subsection{Worked Example: Exact and Approximate Forms}
\label{subsec:worked-exact-approximate-constant}
%%%%%%%%%%%%%%%%%%%%%%%%%%%%%%%%%%%%%%%%%%%%%%%%%%%%%%%%%%%%

\begin{workedexamplebox}{Circumference of a Circle}

Let

\[
r=5.
\]

The exact circumference is

\[
C=2\pi(5).
\]

Thus

\[
C=10\pi.
\]

Using

\[
\pi\approx3.14159,
\]

we obtain

\[
C\approx31.4159.
\]

\begin{answerbox}
\[
\boxed{
C=10\pi
\approx31.4159.
}
\]
\end{answerbox}

\end{workedexamplebox}

%%%%%%%%%%%%%%%%%%%%%%%%%%%%%%%%%%%%%%%%%%%%%%%%%%%%%%%%%%%%
\subsection{Constant Tables in Computation}
\label{subsec:constant-tables-computation}
%%%%%%%%%%%%%%%%%%%%%%%%%%%%%%%%%%%%%%%%%%%%%%%%%%%%%%%%%%%%

Computer algebra systems normally store symbolic constants such as

\[
\boxed{
\pi
}
\]

and

\[
\boxed{
e
}
\]

separately from decimal approximations.

This allows exact simplifications such as

\[
\boxed{
\sin\left(\frac{\pi}{2}\right)=1
}
\]

instead of relying entirely on floating-point approximations.

%%%%%%%%%%%%%%%%%%%%%%%%%%%%%%%%%%%%%%%%%%%%%%%%%%%%%%%%%%%%
\subsection{Floating-Point Constants}
\label{subsec:floating-point-constants}
%%%%%%%%%%%%%%%%%%%%%%%%%%%%%%%%%%%%%%%%%%%%%%%%%%%%%%%%%%%%

A computer approximation of \(\pi\) has finite precision.

Thus internally,

\[
\boxed{
\pi_{\mathrm{machine}}
\neq
\pi
}
\]

as exact mathematical objects.

The machine value is chosen sufficiently close for the intended numerical
precision.

%%%%%%%%%%%%%%%%%%%%%%%%%%%%%%%%%%%%%%%%%%%%%%%%%%%%%%%%%%%%
\subsection{Common Errors}
\label{subsec:mathematical-constants-common-errors}
%%%%%%%%%%%%%%%%%%%%%%%%%%%%%%%%%%%%%%%%%%%%%%%%%%%%%%%%%%%%

\subsubsection{Replacing Exact Constants Too Early}

Replacing

\[
\pi
\]

by

\[
3.14
\]

early in a derivation introduces rounding error.

Keep exact constants symbolic as long as possible.

%%%%%%%%%%%%%%%%%%%%%%%%%%%%%%%%%%%%%%%%%%%%%%%%%%%%%%%%%%%%
\subsubsection{Treating an Approximation as Equality}

Write

\[
\boxed{
\pi
\approx
3.14159,
}
\]

not

\[
\pi
=
3.14159.
\]

The decimal is finite while \(\pi\) is not.

%%%%%%%%%%%%%%%%%%%%%%%%%%%%%%%%%%%%%%%%%%%%%%%%%%%%%%%%%%%%
\subsubsection{Confusing \(e\) With a Variable}

In many contexts,

\[
e
\]

means Euler's number.

However, a text could also define \(e\) as another variable.

Context determines meaning.

%%%%%%%%%%%%%%%%%%%%%%%%%%%%%%%%%%%%%%%%%%%%%%%%%%%%%%%%%%%%
\subsubsection{Confusing Euler's Number and Euler--Mascheroni Constant}

Remember

\[
\boxed{
e
\approx
2.71828
}
\]

while

\[
\boxed{
\gamma
\approx
0.57722.
}
\]

They are unrelated constants despite both being associated historically with
Euler.

%%%%%%%%%%%%%%%%%%%%%%%%%%%%%%%%%%%%%%%%%%%%%%%%%%%%%%%%%%%%
\subsubsection{Confusing \(\pi\) With the Prime-Counting Function}

In number theory,

\[
\boxed{
\pi(x)
}
\]

often denotes the number of primes less than or equal to \(x\).

This is a function, not multiplication by the circle constant.

%%%%%%%%%%%%%%%%%%%%%%%%%%%%%%%%%%%%%%%%%%%%%%%%%%%%%%%%%%%%
\subsubsection{Confusing \(\phi\) With Euler's Totient Function}

The golden ratio is

\[
\boxed{
\phi
=
\frac{1+\sqrt5}{2}.
}
\]

Euler's totient function is often

\[
\boxed{
\phi(n)
}
\]

or

\[
\boxed{
\varphi(n).
}
\]

Context resolves the ambiguity.

%%%%%%%%%%%%%%%%%%%%%%%%%%%%%%%%%%%%%%%%%%%%%%%%%%%%%%%%%%%%
\subsubsection{Assuming Every Named Constant Is Transcendental}

Some are algebraic.

For example,

\[
\boxed{
\sqrt2
}
\]

and

\[
\boxed{
\phi
}
\]

are algebraic.

%%%%%%%%%%%%%%%%%%%%%%%%%%%%%%%%%%%%%%%%%%%%%%%%%%%%%%%%%%%%
\subsubsection{Assuming Every Named Constant Has Known Arithmetic Status}

For some constants, irrationality or transcendence remains unresolved.

Do not infer a classification merely because a constant has a name.

%%%%%%%%%%%%%%%%%%%%%%%%%%%%%%%%%%%%%%%%%%%%%%%%%%%%%%%%%%%%
\subsubsection{Confusing Universal Constants With Model Parameters}

A model may call

\[
\lambda
\]

a constant because it does not vary with time.

That does not make \(\lambda\) a universal mathematical constant.

%%%%%%%%%%%%%%%%%%%%%%%%%%%%%%%%%%%%%%%%%%%%%%%%%%%%%%%%%%%%
\subsubsection{Ignoring Units}

Pure mathematical constants are generally dimensionless.

If a quantity called a constant has physical units, it is usually a physical
or model parameter rather than a pure mathematical constant.

%%%%%%%%%%%%%%%%%%%%%%%%%%%%%%%%%%%%%%%%%%%%%%%%%%%%%%%%%%%%
\subsubsection{Using Excessive Decimal Precision}

If input data are known only to three significant digits, reporting a result
with fifteen decimal places is usually misleading.

Numerical precision should reflect data and model precision.

%%%%%%%%%%%%%%%%%%%%%%%%%%%%%%%%%%%%%%%%%%%%%%%%%%%%%%%%%%%%
\subsection{Quick Decimal Reference}
\label{subsec:quick-decimal-constants}
%%%%%%%%%%%%%%%%%%%%%%%%%%%%%%%%%%%%%%%%%%%%%%%%%%%%%%%%%%%%

\begin{center}
\begin{tabular}{c|l}
\textbf{Constant} & \textbf{Decimal Approximation} \\ \hline
\(\pi\) & \(3.141592653589793\) \\
\(e\) & \(2.718281828459045\) \\
\(\sqrt2\) & \(1.414213562373095\) \\
\(\sqrt3\) & \(1.732050807568877\) \\
\(\phi\) & \(1.618033988749895\) \\
\(\gamma\) & \(0.577215664901533\) \\
\(\ln2\) & \(0.693147180559945\) \\
\(\zeta(2)\) & \(1.644934066848226\) \\
\(\zeta(3)\) & \(1.202056903159594\) \\
\(G\) & \(0.915965594177219\) \\
\(\Omega\) & \(0.567143290409784\)
\end{tabular}
\end{center}

%%%%%%%%%%%%%%%%%%%%%%%%%%%%%%%%%%%%%%%%%%%%%%%%%%%%%%%%%%%%
\subsection{Quick Exact-Form Reference}
\label{subsec:quick-exact-constant-reference}
%%%%%%%%%%%%%%%%%%%%%%%%%%%%%%%%%%%%%%%%%%%%%%%%%%%%%%%%%%%%

\[
\boxed{
\phi
=
\frac{1+\sqrt5}{2}
}
\]

\[
\boxed{
\zeta(2)
=
\frac{\pi^2}{6}
}
\]

\[
\boxed{
\zeta(4)
=
\frac{\pi^4}{90}
}
\]

\[
\boxed{
e
=
\sum_{n=0}^{\infty}
\frac{1}{n!}
}
\]

\[
\boxed{
\gamma
=
\lim_{n\to\infty}
\left(
H_n-\ln n
\right)
}
\]

\[
\boxed{
G
=
\sum_{n=0}^{\infty}
\frac{(-1)^n}{(2n+1)^2}
}
\]

\[
\boxed{
\Omega e^\Omega=1.
}
\]

%%%%%%%%%%%%%%%%%%%%%%%%%%%%%%%%%%%%%%%%%%%%%%%%%%%%%%%%%%%%
\subsection{Exercises}
\label{subsec:mathematical-constants-exercises}
%%%%%%%%%%%%%%%%%%%%%%%%%%%%%%%%%%%%%%%%%%%%%%%%%%%%%%%%%%%%

\begin{exercise}[1]
State the approximate value of \(\pi\) to four decimal places.
\end{exercise}

\begin{exercise}[1]
State the approximate value of \(e\) to four decimal places.
\end{exercise}

\begin{exercise}[1]
Define \(\pi\) geometrically.
\end{exercise}

\begin{exercise}[1]
Give one limit definition of \(e\).
\end{exercise}

\begin{exercise}[1]
Define the golden ratio.
\end{exercise}

\begin{exercise}[1]
Define the Euler--Mascheroni constant.
\end{exercise}

\begin{exercise}[2]
Compute the circumference of a circle of radius \(7\) exactly and
approximately.
\end{exercise}

\begin{exercise}[2]
Compute the area of a circle of radius \(4\) exactly and approximately.
\end{exercise}

\begin{exercise}[2]
Convert \(135^\circ\) to radians.
\end{exercise}

\begin{exercise}[2]
Convert \(5\pi/6\) radians to degrees.
\end{exercise}

\begin{exercise}[2]
Use

\[
\phi
=
\frac{1+\sqrt5}{2}
\]

to verify

\[
\phi^2=\phi+1.
\]
\end{exercise}

\begin{exercise}[2]
Verify

\[
\frac1\phi
=
\phi-1.
\]
\end{exercise}

\begin{exercise}[2]
Compute the first six powers of \(i\).
\end{exercise}

\begin{exercise}[2]
Use Euler's formula to evaluate

\[
e^{i\pi}.
\]
\end{exercise}

\begin{exercise}[2]
Explain why \(\sqrt2\) is algebraic.
\end{exercise}

\begin{exercise}[2]
State whether \(\pi\) is algebraic or transcendental.
\end{exercise}

\begin{exercise}[2]
State whether \(e\) is algebraic or transcendental.
\end{exercise}

\begin{exercise}[2]
State whether \(\phi\) is algebraic or transcendental.
\end{exercise}

\begin{exercise}[2]
Compute

\[
\zeta(2)
\]

using its exact formula.
\end{exercise}

\begin{exercise}[2]
Approximate

\[
\ln2
\]

using the first several terms of the alternating harmonic series.
\end{exercise}

\begin{exercise}[3]
Approximate \(e\) using

\[
\sum_{n=0}^{N}
\frac1{n!}.
\]
\end{exercise}

\begin{exercise}[3]
Approximate \(\pi\) using the Gregory--Leibniz series and comment on its
convergence speed.
\end{exercise}

\begin{exercise}[3]
Verify numerically that

\[
H_n-\ln n
\]

approaches \(\gamma\).
\end{exercise}

\begin{exercise}[3]
Use Stirling's approximation to estimate \(20!\).
\end{exercise}

\begin{exercise}[3]
Compare the Stirling approximation of \(20!\) with its exact value.
\end{exercise}

\begin{exercise}[3]
Explain why the constant

\[
1/\sqrt{2\pi}
\]

appears in the standard normal density.
\end{exercise}

\begin{exercise}[3]
Show that the golden ratio arises as the positive root of

\[
x^2-x-1=0.
\]
\end{exercise}

\begin{exercise}[3]
Show that ratios of consecutive Fibonacci numbers approach the positive root
of

\[
x^2=x+1.
\]
\end{exercise}

\begin{exercise}[3]
Explain the difference between a constant of integration and a universal
mathematical constant.
\end{exercise}

\begin{exercise}[3]
Explain why a finite decimal approximation to \(\pi\) is rational even
though \(\pi\) is irrational.
\end{exercise}

\begin{exercise}[3]
Explain why mathematical constants such as \(\pi\) are dimensionless.
\end{exercise}

\begin{exercise}[3]
Derive

\[
\int_0^1
\frac{1}{1+x^2}\,dx
=
\frac{\pi}{4}.
\]
\end{exercise}

\begin{exercise}[3]
Investigate Wallis' product numerically.
\end{exercise}

\begin{exercise}[4]
Derive the Gaussian integral

\[
\int_{-\infty}^{\infty}
e^{-x^2}\,dx
=
\sqrt{\pi}.
\]
\end{exercise}

\begin{exercise}[4]
Use the gamma function to prove

\[
\Gamma\left(\frac12\right)
=
\sqrt{\pi}.
\]
\end{exercise}

\begin{exercise}[4]
Derive the asymptotic relation

\[
H_n
=
\ln n+\gamma+o(1).
\]
\end{exercise}

\begin{exercise}[4]
Study Euler's solution of the Basel problem conceptually.
\end{exercise}

\begin{exercise}[4]
Derive Binet's formula for Fibonacci numbers.
\end{exercise}

\begin{exercise}[4]
Investigate the convergence of the ratio

\[
\frac{F_{n+1}}{F_n}
\]

to \(\phi\).
\end{exercise}

\begin{exercise}[4]
Study the appearance of the Feigenbaum constant in the logistic map.
\end{exercise}

\begin{exercise}[4]
Investigate Catalan's constant through its series and integral
representations.
\end{exercise}

\begin{exercise}[4]
Study the Lambert \(W\) relation

\[
\Omega=W(1).
\]
\end{exercise}

\begin{exercise}[4]
Compare numerical algorithms for computing \(\pi\).
\end{exercise}

\begin{exercise}[4]
Compare numerical algorithms for computing \(e\).
\end{exercise}

\begin{exercise}[4]
Investigate how floating-point precision affects stored approximations of
mathematical constants.
\end{exercise}

%%%%%%%%%%%%%%%%%%%%%%%%%%%%%%%%%%%%%%%%%%%%%%%%%%%%%%%%%%%%
\subsection{Conceptual Summary}
\label{subsec:mathematical-constants-summary}
%%%%%%%%%%%%%%%%%%%%%%%%%%%%%%%%%%%%%%%%%%%%%%%%%%%%%%%%%%%%

Mathematical constants are fixed values that arise naturally from
mathematical structure.

The most fundamental examples include

\[
\boxed{
\pi,
\qquad
e,
\qquad
i,
\qquad
\phi,
\qquad
\gamma.
}
\]

The circle constant

\[
\boxed{
\pi
}
\]

connects geometry, trigonometry, analysis, probability, and complex
mathematics.

Euler's number

\[
\boxed{
e
}
\]

governs exponential growth, logarithms, differential equations, probability,
and asymptotic analysis.

The imaginary unit

\[
\boxed{
i
}
\]

extends real arithmetic to the complex numbers.

The golden ratio

\[
\boxed{
\phi
=
\frac{1+\sqrt5}{2}
}
\]

connects algebra, geometry, and recurrence relations.

The Euler--Mascheroni constant

\[
\boxed{
\gamma
=
\lim_{n\to\infty}
(H_n-\ln n)
}
\]

arises in asymptotic analysis and number theory.

Other constants such as

\[
\boxed{
\zeta(3),
\quad
G,
\quad
\Omega,
}
\]

and the Feigenbaum constants illustrate how fixed numerical values emerge
from infinite series, integrals, limits, nonlinear dynamics, and special
functions.

The central computational principle is

\[
\boxed{
\text{preserve exact constants symbolically until approximation is needed}.
}
\]

%%%%%%%%%%%%%%%%%%%%%%%%%%%%%%%%%%%%%%%%%%%%%%%%%%%%%%%%%%%%
\subsection{Section Connections}
\label{subsec:mathematical-constants-connections}
%%%%%%%%%%%%%%%%%%%%%%%%%%%%%%%%%%%%%%%%%%%%%%%%%%%%%%%%%%%%

\chapterconnection{
Mathematical Constants collects fixed numerical values that recur throughout
the textbook. The constants \(\pi\) and \(e\) appear throughout Geometry,
Analysis, Differential Equations, Probability, and Mathematical Finance;
\(i\) is fundamental to Complex Analysis and spectral methods; \(\phi\)
connects Algebra, Geometry, and recurrence relations; \(\gamma\), zeta
values, and related constants connect Analysis and Number Theory; and
universal scaling constants arise in Dynamical Systems. The next reference
section, Algebraic Identities, collects standard factoring formulas,
polynomial identities, exponent rules, logarithmic identities, radical
relations, and other algebraic formulas used throughout the text.
}

%%%%%%%%%%%%%%%%%%%%%%%%%%%%%%%%%%%%%%%%%%%%%%%%%%%%%%%%%%%%
% SECTION SOURCE TRACKING
%%%%%%%%%%%%%%%%%%%%%%%%%%%%%%%%%%%%%%%%%%%%%%%%%%%%%%%%%%%%

%------------------------------------------------------------
% 14.4 Mathematical Constants
%------------------------------------------------------------
%
% Source basis:
%   - Standard mathematical constants and identities.
%   - Standard analysis, geometry, number theory, probability,
%     complex analysis, asymptotic, and dynamical-systems formulas.
%
% Used for:
%   - pi;
%   - Euler's number;
%   - imaginary unit;
%   - golden ratio;
%   - Euler--Mascheroni constant;
%   - square-root constants;
%   - natural logarithm of two;
%   - zeta values;
%   - Apery's constant;
%   - Catalan's constant;
%   - Feigenbaum constant;
%   - Khinchin's constant;
%   - Glaisher--Kinkelin constant;
%   - omega constant;
%   - plastic constant;
%   - geometric formulas;
%   - exponential and logarithmic formulas;
%   - Euler's identity;
%   - Fibonacci relations;
%   - harmonic-number asymptotics;
%   - Gaussian normalization;
%   - Fourier constants;
%   - differential-equation constants;
%   - combinatorial appearances of e;
%   - Stirling's approximation;
%   - constants of integration;
%   - algebraic and transcendental constants;
%   - series, integral, product, and limit representations;
%   - exact versus approximate calculation;
%   - computational precision;
%   - applications and exercises.
%
% Notes:
%   - Algebraic Identities is developed next in Section 14.5.
%   - Geometric Formulas follows in Section 14.6.
%
%%%%%%%%%%%%%%%%%%%%%%%%%%%%%%%%%%%%%%%%%%%%%%%%%%%%%%%%%%%%